\subsection{The dyadic sixteen-torsion problem reduces to two supersingular curves}

\begin{lemma}[Finite dyadic inertia and integral elliptic conjugacy]
\label{lem:B-RD44-conjugacy}
\status{Proved; complete independent parent and obstructions/crystals proof, source and exact-code audits passed.}
Let $k$ be algebraically closed of characteristic two and let
$G_0=E_0[2^\infty]$, where $E_0:y^2+y=x^3$.
Every finite subgroup of $\operatorname{Aut}(G_0)$ is conjugate,
by an element of $\operatorname{Aut}(G_0)$ itself, into
$\operatorname{Aut}(E_0)$. The possible groups are
\[
 1,\ C_2,\ C_3,\ C_4,\ C_6,\ Q_8,\ \operatorname{SL}_2(\mathbf F_3).
 \tag{B-RD44.1}
\]
Their reduced norms are all one.
\end{lemma}
\begin{proof}
The height-two Dieudonne calculation of Lemma~\ref{lem:B-RD31-finite},
with $3$ replaced by $2$, identifies the rational endomorphism algebra
with the quaternion division algebra $D/\mathbf Q_2$, and
$\operatorname{End}(G_0)$ with its valuation ring $\mathcal O_D$.
Its residue field is $\mathbf F_4$ and the reduction kernel on units
is pro-$2$. Normalize $v_D$ by $v_D(2)=2$.
For a finite subgroup $H$, its reduction kernel $P$ is consequently
a normal Sylow two-subgroup, and $H/P$ has order one or three.

Every element of $D$ satisfies a quadratic polynomial over $\mathbf Q_2$.
Thus no element of $P$ has order eight, since
$[\mathbf Q_2(\zeta_8):\mathbf Q_2]=4$; every element has order at
most four. The only involution in a division ring of characteristic
zero is $-1$, by $(x-1)(x+1)=0$.
If two order-four elements $i,j$ generate different cyclic subgroups,
then $ij$ is neither $1$ nor $-1$, so $(ij)^2=-1$ and $ij=-ji$.
They generate $Q_8$. Any further order-four element outside this
quaternion group would anticommute with both $i,j$. In the basis
$1,i,j,ij$ its simultaneous anticommutant is $\mathbf Q_2ij$;
squaring then forces the element to be $ij$ or $-ij$, a contradiction.
This proves $P=1,C_2,C_4$, or $Q_8$.

If $H/P$ has order three, Cauchy's theorem gives an element $z$ of
order three and $H=P\rtimes\langle z\rangle$.
For $P=1,C_2$ this gives $C_3,C_6$.
For $P=C_4$ the action is trivial, since
$\operatorname{Aut}(C_4)=C_2$; this would give an element of order
twelve, impossible because $\mathbf Q_2(\zeta_{12})$ has degree four.
For $P=Q_8$ the action must be nontrivial for the same reason.
It cyclically permutes the three order-four subgroups, giving the
binary tetrahedral group $Q_8\rtimes C_3=\operatorname{SL}_2(\mathbf F_3)$.
This recovers the degree-two case of Hewett's general classification
\cite[Theorem~1.4 and Corollary~1.5, p.~520]{Hewett1995Division}.
The argument above proves the particular classification used here
without importing a conjugacy assertion from that theorem.

We now prove integral conjugacy. Choose $\omega\in\mathbf F_4$
with $\omega^2+\omega+1=0$. The maps
\[
 (x,y)\longmapsto
 (u^2x+s^2,\ y+u^2sx+t),\qquad
 u\in\mathbf F_4^\times,\ s\in\mathbf F_4,\ t^2+t=s^3,
 \tag{B-RD44.2}
\]
are twenty-four distinct automorphisms fixing the origin.
Substitution verifies their equations and closure under composition.
The eight with $u=1$ form $Q_8$: the six with $s\ne0$ square
to $(x,y)\mapsto(x,y+1)=[-1]$.
Choose generators $i,j$ of this quaternion subgroup.
Its rational span is $D$, with $i^2=j^2=-1$ and $ij=-ji$.
The eight order-three automorphisms in (B-RD44.2) induce the eight
order-three automorphisms of $Q_8$; this follows either by composing
the displayed maps, or by choosing cyclic images of the three axes
and the two independent signs. Denote this group of twenty-four
automorphisms by $H_0$.

An order-four element $a$ can be conjugated to $i$ by $a+i$,
unless $a=-i$, which already generates the desired subgroup:
$(a+i)a=i(a+i)$.
The centralizer of $i$ contains $1+i$, of norm two and valuation one.
Multiplication of this conjugator on the left by any power of $1+i$
does not change its conjugation of $a$, and makes its valuation zero.
Thus the cyclic order-four conjugacy is integral.

Put $\zeta=(-1+i+j+ij)/2$.
Its conjugation cyclically permutes the quaternion axes and has
order three. Choose the order-three member of $H_0$ inducing that
same automorphism of $Q_8$. Their quotient is central in $D$,
and its cube is one, so it is one; thus $\zeta\in H_0$.
For any order-three $z$, the element $d=z+\zeta+1$ satisfies
$dz=\zeta d$. If $d=0$, then $z=\zeta^{-1}$ and its cyclic
subgroup is already the required one. Otherwise $d$ is invertible.
The element $i-j$ has norm two, valuation one, and conjugates
$\zeta$ to $\zeta^{-1}$. Hence multiplying $d$ on the left by
a suitable integral power of $i-j$ makes it a unit while still
conjugating $\langle z\rangle$ to $\langle\zeta\rangle$.
This also treats $C_6$, since the involution is $-1$.

For a quaternion subgroup, first put one generator equal to $i$
by the preceding integral conjugacy. Its other generator is $cj$,
where $c\in\mathbf Q_2(i)$ and $c\bar c=1$.
If $c\ne-1$, take $b=1+c$; if $c=-1$, take $b=i$.
In both cases $b/\bar b=c$, so $b^{-1}(cj)b=j$ while $b$ commutes
with $i$. The element $1+i$ normalizes the target $Q_8$, sending
$j$ to $ij$, and again has valuation one. Left multiplication
of $b^{-1}$ by a suitable power of $1+i$ therefore makes the
conjugator integral without changing the target subgroup.

Finally put the quaternion subgroup of an order-twenty-four $H$
inside $H_0$ by this integral conjugacy. An element $z\in H$ of
order three induces the same automorphism of $Q_8$ as some
order-three $z_0\in H_0$. Then $zz_0^{-1}$ centralizes the
rational span $D$ of $Q_8$, so is a scalar in $\mathbf Q_2^\times$.
Its cube is one, forcing it to be one. Hence $H\subset H_0$.
The trivial and order-two cases require no conjugacy.
All correcting elements have valuation zero at the end and
therefore belong to $\mathcal O_D^\times$, as required.
Finally order-four generators have minimal polynomial $T^2+1$,
and order-three generators have $T^2+T+1$; their norms, and that
of $-1$, are one. They generate every group in (B-RD44.1).
\end{proof}

\begin{proposition}[Exact dyadic elliptic realization]
\label{prop:B-RD44-realization}
\status{Proved; complete independent parent and obstructions/crystals proof, source and exact-code audits passed.}
Let $K$ be a complete discretely valued field of mixed characteristic
$(0,2)$ with algebraically closed residue field. Every height-two
$2$-divisible group $M/K$ which has potentially good one-dimensional
supersingular reduction is isomorphic to $E[2^\infty]$ for an
elliptic curve $E/K$. The isomorphism preserves the full integral
Tate lattice and the labelled finite special-inertia action.
\end{proposition}
\begin{proof}
Choose a finite Galois good extension $L/K$ and the good model
$\mathcal G/\mathcal O_L$. Tate's integral full faithfulness
extends its generic semilinear descent isomorphisms and their
inverses uniquely, with their cocycle relations
\cite[Theorem~4 and Corollaries~1--2, p.~181]{Tate1967PDivisible}.
Lazard's height classification
\cite[Theorem~IV and proof, pp.~270--271]{Lazard1955Formal}
identifies the special formal group with $G_0$.
Lemma~\ref{lem:B-RD44-conjugacy} makes the entire labelled finite
action an action by actual automorphisms of $E_0$.

The labelled Serre--Tate construction and algebraization in the
proof of Proposition~\ref{prop:B-RD35} now apply with $2$ in place
of $3$: their deformation bases are $\mathcal O_L/\pi^n$, on
which $2$ is nilpotent, and their precise inputs impose no
odd-prime restriction. Namely Katz's Theorem~1.2.1
\cite[p.~143 and proof pp.~143--146]{Katz1981SerreTate}
produces compatible elliptic deformations with the chosen BT
isomorphisms and semilinear action. The compatible ample line
bundles $\mathcal O(3\cdot0)$ and EGA~III$_1$, Theorems~5.4.1
and~5.4.5, algebraize the elliptic scheme and its operations
\cite[pp.~156--158]{EGAIII1}.
For each $r$, completeness of the finite free Hopf algebras
algebraizes the isomorphism on $2^r$-torsion, compatibly in $r$.
Serre--Tate uniqueness and proper formal full faithfulness lift
and algebraize the labelled descent datum and all cocycles.
Effective generic polarized Galois descent yields $E/K$ and
$E[2^\infty]\simeq M$. As in B-RD35, only the generic elliptic
curve is descended across the ramified extension; its good
integral model is not asserted to descend.
\end{proof}

\begin{lemma}[No rational sixteen-point in the unramified supersingular case]
\label{lem:B-RD44-point}
\status{Proved; complete independent parent and obstructions/crystals proof, source and exact-code audits passed.}
Let $K$ have perfect residue field of characteristic two and
absolute ramification index one. If an elliptic curve $E/K$ has
potentially supersingular good reduction, then
$E(K)[2^\infty]$ has exponent at most eight.
In particular it has no point of exact order sixteen.
\end{lemma}
\begin{proof}
Pass to the completed maximal unramified extension if necessary;
this preserves the absolute ramification and any asserted torsion
point. Choose a minimal integral Weierstrass equation.
Its coefficient $a_1$ lies in $2\mathcal O_K$.
For good supersingular reduction this is the vanishing of the
characteristic-two Hasse coefficient. For bad reduction, potential
goodness gives $v(j)\ge0$ and $v(\Delta)>0$, whence $v(c_4)>0$;
the identity $c_4\equiv a_1^4\pmod2$ proves the same assertion.

The formal parameter $T=-x/y$ identifies the reduction kernel
$E_1(K)$ with the maximal ideal and has integral multiplication
series
\[
             [2](T)=2T-a_1T^2+\sum_{j\ge3}b_jT^j.
 \tag{B-RD44.3}
\]
These integral formal-group assertions are Tate's Theorems~4.1--4.2
and the expansions in his Section~4
\cite[pp.~41--44]{Tate1975Algorithm}.
If $v(T)\ge1$, the first term of (B-RD44.3) has value $1+v(T)$,
strictly less than $1+2v(T)$ and than every $jv(T)$ for $j\ge3$.
Consequently $[2]T\ne0$ for every nonzero formal point, and
$E_1(K)[2^\infty]=0$.

If reduction is good supersingular, its special curve has no
nonidentity geometric two-primary torsion. Every rational
two-primary point lies in $E_1(K)$ and is therefore zero.
Otherwise reduction is additive: potential good reduction excludes
the multiplicative case. The smooth connected reduction group is
$\mathbf G_a$, so $E_0(K)/E_1(K)$ has exponent two. The geometric
component group's two-primary subgroup has exponent at most four.
Both assertions hold in residue characteristic two: in Tate's
Section~6 table, p.~46, these are the component-group and
identity-component rows above the explicit boundary marking the
later assumptions $\operatorname{char}k\ne2,3$
\cite[Section~6, pp.~45--46]{Tate1975Algorithm}.
Thus any rational two-primary point $P$ has $[4]P\in E_0(K)$
and $[8]P\in E_1(K)$. The latter group is two-primary torsion-free,
so $[8]P=0$.
\end{proof}

\begin{proposition}[B-RD44: the mixed case is excluded by sixteen-torsion]
\label{prop:B-RD44}
\status{Proved; complete independent parent and obstructions/crystals proof, source and exact-code audits passed.}
Suppose $p=2$, $e_K=1$, and $M,N$ have height two and satisfy
the manuscript monodromy condition. If
$M[16]\simeq N[16]$ but their prime-to-two special-inertia
characters differ, then both groups themselves have potentially
good supersingular reduction. After completed unramified base
change they are the exact BT groups of two elliptic curves.
Their faithful finite inertia images have respective orders
\[
       m_1\in\{1,2,4,8\},\qquad m_2\in\{3,6,24\},
 \tag{B-RD44.4}
\]
in some order; their joint good-reduction degree belongs to
$\{3,6,12,24,48,96,192\}$.
No assertion of sixteen-torsion sufficiency for these remaining
pairs is made.
\end{proposition}
\begin{proof}
Work after completed maximal unramified base change. Write $d_M,d_N$
for the sums of the dimensions of the good graded models.
They belong to $\{0,1,2\}$. The determinant characters are
$\chi_2^{d_M}\epsilon_M,\chi_2^{d_N}\epsilon_N$, where
$\epsilon_M,\epsilon_N$ are finite quadratic characters, by the
height-one determinant classification already used in B-RD18 and
B-RD31. The cyclotomic image modulo sixteen is all
$(\mathbf Z/16)^\times$: the polynomial $\Phi_{16}(1+T)$ is
Eisenstein over $\mathcal O_K$ because $v_K(2)=1$, so adjoining
$\zeta_{16}$ still has degree eight, also for a transcendental
perfect residue field.
Equality of determinants modulo sixteen forces $d_M=d_N$:
if their difference is $\pm1$, its cyclotomic character has
order four and cannot be a sign character; if it is $\pm2$,
an inertia element with cyclotomic value $3$ has squared value
$9$, which is neither $1$ nor $-1$ modulo sixteen.

For common dimension zero, all good graded groups are etale,
and the prime-to-two Brauer character is read directly from the
semisimplified two-torsion. Dimension two reduces to this by
Cartier duality. For common dimension one, each group either
has a filtration by height-one pieces, or is itself potentially
good supersingular. A potentially ordinary good graded piece
has its canonical connected--etale filtration, which descends
uniquely and gives the first alternative.
Height-one finite characters have values in
$(\mathbf Q_2^\times)_{\mathrm{tors}}=\{1,-1\}$; hence every
height-one-filtered group has prime-to-two special character two.

It remains to exclude a comparison between a potentially
supersingular $M$ and a height-one-filtered $N$.
Proposition~\ref{prop:B-RD44-realization} gives $M=E[2^\infty]$
with its exact Tate lattice. The first height-one subgroup of
$N[16]$ transports to a cyclic direct summand $C\subset E[16]$
with character $\chi_2^b\epsilon$, where $b=0$ or $1$ and
$\epsilon$ is quadratic. If $b=0$, the quadratic twist by
$\epsilon$ has a rational point of order sixteen.
If $b=1$, use the elliptic quotient $E'=E/C$.
The image of $E[16]$ is the cyclic subgroup $E[16]/C$ of
$E'[16]$; the Weil pairing, or the determinant $\det E[16]=\chi_2$,
gives this quotient the character $\epsilon^{-1}$.
The corresponding quadratic twist of $E'$ again has a rational
point of order sixteen. Isogenies and quadratic twists preserve
potentially supersingular reduction. Both conclusions contradict
Lemma~\ref{lem:B-RD44-point}.

Both remaining groups are therefore potentially supersingular,
and Proposition~\ref{prop:B-RD44-realization} applies to each.
In the list (B-RD44.1), every odd-order element is either the
identity, of trace two, or has order three, of trace minus one.
If both faithful inertia orders are divisible by three, their
characters agree as functions on prime-to-two inertia: the tame
procyclic quotient has a unique quotient of order three, so the
two nontrivial cubic quotients have the same kernel. If neither
order is divisible by three, both characters are two. Different
characters therefore require exactly one order divisible by three,
which gives (B-RD44.4). The joint wild image embeds in the product
of two groups of orders at most eight, and its tame quotient has
order three. Its order is consequently $3\cdot2^r$, $0\le r\le6$.
The faithful-image descent of Lemma~\ref{lem:B-RD17-degree}
makes this a simultaneous good-reduction extension.
\end{proof}

\begin{remark}[Scope and finite verification]
Proposition~\ref{prop:B-RD43} proves that eight-torsion can fail
at $(p,e,h)=(2,1,2)$. The preceding result removes the mixed
case at the next level, but leaves the all-supersingular
sixteen-torsion comparison open. No finite elimination has been
executed for those pairs.
The certificate \path{code/verify_p2_height_two_mixed_gate.py}
checks the twenty-four special automorphisms, their group orders
and quaternion relations, the explicit conjugator and normalizer
identities, the formal valuation inequalities, and the determinant
dimension separation modulo sixteen. Its exact scope and runtime
are retained in the adjacent JSON output. The author run passed
in $0.07470255399675807$ seconds; these computations do
not decide the remaining pairs.
\end{remark}

\paragraph{Independent verification.}
The root rerun took
$0.06556921700394014$ seconds (actual audit record
2026-11-19 at 13:57:27.954980 UTC); the obstructions/crystals
agent reran it in $0.06939755400526337$ seconds.
Both independently read the full proof and the specified original
sources. The author output was preserved. These assertions exclude
the mixed sixteen-torsion comparison and give an exact elliptic
realization; they do not decide the remaining comparison between
two potentially supersingular objects.
